% ============================================================
% UNIDAD 9: HERRAMIENTAS Y FERRETERÍA APLICADA A LA AVIÓNICA
% ============================================================
\clearpage
\section*{UNIDAD 9: HERRAMIENTAS Y FERRETERÍA}
\subsection*{Herramental del taller y materiales eléctricos y electrónicos}

\begin{recuadro}
\textbf{Objetivo de la unidad:} conocer, seleccionar y usar correctamente las herramientas e
instrumentos del taller de aviónica, y los elementos de ferretería, materiales y consumibles
eléctricos y electrónicos aprobados para aeronaves.
\end{recuadro}

% ------------------------------------------------------------
\section*{1. HERRAMIENTAS DE CABLEADO}
\begin{center}
\small
\begin{tabular}{|p{0.30\linewidth}|p{0.58\linewidth}|}
\hline
\textbf{Herramienta} & \textbf{Uso y cuidados} \\ \hline
Peladoras de precisión & Con cuchillas para cada sección; no usar con cables de otra aislación. \\
\hline
Peladora térmica & Para aislaciones finas y cables de calibre pequeño. \\ \hline
Pinzas de crimpar con trinquete & Terminales y empalmes; matriz según color. Calibración
periódica. \\ \hline
Crimpadoras M22520/1-01 y /2-01 & Contactos de conectores, con posicionadores y torretas. \\
\hline
Herramientas de inserción y extracción & Una por tamaño de contacto; se reemplazan si se doblan.
\\ \hline
Alicates de corte al ras & Cortes limpios de cables y precintos sin dejar filos. \\ \hline
Herramienta tensora de precintos & Aprieta y corta los precintos a la tensión correcta. \\ \hline
Pistola de calor & Termocontraíbles y manguitos de soldadura, con boquillas reflectoras. \\ \hline
Estación de soldado y desoldado & Con temperatura controlada y punta adecuada. \\ \hline
Tablero de mazos & Para armar los mazos a escala. \\ \hline
Etiquetadora o impresora de manguitos & Identificación de cables. \\ \hline
\end{tabular}
\end{center}

% ------------------------------------------------------------
\section*{2. HERRAMIENTAS MECÁNICAS}
\begin{itemize}
  \item \textbf{Destornilladores} de punta correcta (Phillips, Torx, plana) y en buen estado: una
        punta gastada daña la cabeza del tornillo y el panel.
  \item \textbf{Llaves} de boca, estriadas y tubos en pulgadas.
  \item \textbf{Llaves dinamométricas} (torquímetros) para tornillos, tuercas y conectores, con
        calibración vigente.
  \item \textbf{Pinza para alambre de frenar} y alambre de frenar del diámetro y material
        indicados.
  \item \textbf{Sacabocados de panel} (\textit{panel punch}) para los orificios de 2 1/4 y 3 1/8
        pulgadas de los instrumentos; mechas, avellanadores y limas.
  \item \textbf{Remachadora} y herramientas para tuercas remachables y tuercas con placa
        (\textit{nutplates}).
  \item \textbf{Espejo de inspección, linterna y lupa.}
\end{itemize}

% ------------------------------------------------------------
\section*{3. INSTRUMENTOS DE MEDICIÓN}
\begin{center}
\small
\begin{tabular}{|p{0.30\linewidth}|p{0.58\linewidth}|}
\hline
\textbf{Instrumento} & \textbf{Uso} \\ \hline
Multímetro digital & Tensión, corriente, resistencia, continuidad. \\ \hline
Miliohmímetro & Resistencia de masas y \textit{bonding}. \\ \hline
Medidor de aislamiento & Aislamiento de cables, con la tensión adecuada y \textbf{siempre con los
equipos desconectados}. \\ \hline
Osciloscopio & Señales de datos (ARINC 429, CAN, RS-232), sensores, ruido. \\ \hline
Banco de prueba pitot-estático & Ensayo de pérdidas y verificación de altímetros, velocímetros y
ADC (Unidad 11). \\ \hline
Simuladores de sensores & Década de resistencias, calibrador de termocuplas, generador de pulsos
para simular sensores de motor. \\ \hline
Equipos de prueba de navegación y transpondedor & Generadores de señales de VOR/ILS y de radar
secundario. \\ \hline
\end{tabular}
\end{center}

\paragraph*{Control de herramientas.} En un taller aeronáutico cada herramienta está
\textbf{identificada} y \textbf{controlada}: tableros de sombras o cajas con su lugar marcado, para
que al terminar una tarea se compruebe que no quedó ninguna dentro de la aeronave; y las que
requieren calibración llevan su etiqueta con la fecha de vencimiento.

% ------------------------------------------------------------
\section*{4. FERRETERÍA ELÉCTRICA Y ELECTRÓNICA APLICADA}

\subsubsection*{4.1. Protección de circuitos}
\begin{itemize}
  \item \textbf{Disyuntores (\textit{circuit breakers}):} protegen el \textbf{cable} contra
        sobrecorriente. Se dimensionan según la sección del cable y el consumo del equipo, como
        indica el fabricante. Los de tipo extraíble (\textit{push-pull}) permiten además abrir el
        circuito manualmente.
  \item \textbf{Fusibles:} en algunos circuitos y equipos, con repuestos a bordo según la norma.
  \item \textbf{Limitadores de corriente:} protecciones de barras de alta corriente.
\end{itemize}

\subsubsection*{4.2. Conmutación y control}
\begin{itemize}
  \item \textbf{Interruptores} de palanca y de presión con norma aeronáutica, del amperaje correcto.
  \item \textbf{Relés y contactores:} con diodo de supresión en paralelo con la bobina cuando lo
        indica el plano, para evitar los picos de tensión al desconectarla.
  \item \textbf{Interruptor maestro de aviónica:} desconecta los equipos durante el arranque del
        motor para protegerlos de los transitorios.
\end{itemize}

\subsubsection*{4.3. Ferretería de montaje}
Los tornillos, tuercas y arandelas de una aeronave se identifican por su \textbf{norma (AN, MS,
NAS)}, no por su medida: por ejemplo, tornillos de cabeza fresada MS24694 o de cabeza alomada
MS35206, tuercas autofrenantes AN365 o MS21042, arandelas AN960 y tuercas con placa para fijar
paneles e instrumentos. Se usan siempre los que indica el IPC o el manual de instalación; nunca
ferretería comercial sin identificación.

\subsubsection*{4.4. Materiales y consumibles}
\begin{itemize}
  \item Cables aprobados (Unidad 6), terminales, contactos, empalmes y manguitos.
  \item Termocontraíbles (norma AMS-DTL-23053), manguitos de protección, espirales y fundas
        trenzadas.
  \item Hilo de atar, precintos, abrazaderas acolchadas, pasacables y protectores de borde.
  \item Estaño, fundente y limpiadores aprobados.
  \item Selladores, compuestos para juntas y pastas conductoras, según el manual.
  \item Etiquetas y placas de identificación de los disyuntores e interruptores.
\end{itemize}

\begin{recuadro}
\textbf{Trazabilidad.} Todo material que se instala en una aeronave debe poder rastrearse: se
compra a proveedores confiables con \textbf{certificado de conformidad}, se guarda identificado y,
en el caso de los equipos, con su certificado de liberación (por ejemplo, el 8130-3). Un terminal
o un cable ``sin papeles'' puede ser una \textbf{parte no aprobada} (o falsificada) y no se instala.
\end{recuadro}

% ------------------------------------------------------------
\section*{5. ACTIVIDADES DE LA UNIDAD}
\begin{enumerate}
  \item Hacer el inventario de las herramientas del taller con su función, estado y, si
        corresponde, fecha de calibración.
  \item Armar un tablero de sombras para las herramientas de cableado.
  \item Dimensionar el disyuntor y el cable para un equipo de 2,5 A en 14 V según su manual de
        instalación.
  \item Identificar por su norma una muestra de tornillos, tuercas y arandelas.
  \item Hacer el pedido de materiales (con números de parte) para instalar un G5 en un panel.
\end{enumerate}
